\section{三角関数で根号を簡単にする}\label{節:三角置換の初歩}

$a>0$ として，$\sqrt{a^2-x^2}$ には平方の差がある。三角関数の関係式 $1-\sin^2\theta=\cos^2\theta$ と形が同じだから，$x=a\sin\theta$ と置くことを考える。すると
\[
 \sqrt{a^2-x^2}
 =a\sqrt{1-\sin^2\theta}
 =a\sqrt{\cos^2\theta}
 =a|\cos\theta|.
\]
最後は $a\cos\theta$ ではなく，まず絶対値付きである。$\sqrt{z^2}=|z|$ だからである。そのうえで，たとえば $-\frac\pi2\le\theta\le\frac\pi2$ を選べば $\cos\theta\ge0$ となり，絶対値を外せる。この区間では $\sin\theta$ が一対一でもあるので，$-a\le x\le a$ の各 $x$ に対して $\theta$ が一つに決まる。

\begin{figure}[htbp]
 \centering
 \begin{tikzpicture}[x=2cm,y=1.5cm]
  \draw[-{Stealth}] (-1.95,0)--(1.95,0) node[right] {$\theta$};
  \draw[-{Stealth}] (0,-1.3)--(0,1.35) node[above] {$\sin\theta$};
  \draw[thick,blue,domain=-1.570796:1.570796,samples=70]
    plot (\x,{sin(\x r)});
  \draw[dashed] (-1.570796,0)--(-1.570796,-1)--(0,-1);
  \draw[dashed] (1.570796,0)--(1.570796,1)--(0,1);
  \node[above] at (-1.570796,0) {$-\frac\pi2$};
  \node[below] at (1.570796,0) {$\frac\pi2$};
  \node[above left] at (0,0) {$0$};
  \node[left] at (0,1) {$1$};
  \node[right] at (0,-1) {$-1$};
 \end{tikzpicture}
 \caption{正弦はこの区間で狭義単調増加し，逆向きの対応を一つに決められる。}
 \label{図:正弦の枝}
\end{figure}

\begin{example}{半円の根号を三角関数へ換える}
$a>0$ とする。$\int_0^a\sqrt{a^2-x^2}\dd x$ を求めよ。
\end{example}

ここでは $x\ge0$ だから，$x=a\sin\theta$ と置き，$0\le\theta\le\frac\pi2$ を選ぶ。$x=0$ は $\theta=0$，$x=a$ は $\theta=\frac\pi2$ に対応する。さらに
\[
 \dd x=a\cos\theta\dd\theta,\qquad
 \sqrt{a^2-x^2}=a\cos\theta
\]
なので，
\[
\begin{aligned}
 \int_0^a\sqrt{a^2-x^2}\dd x
 &=a^2\int_0^{\frac\pi2}\cos^2\theta\dd\theta\\
 &=a^2\left[\frac\theta2+\frac14\sin(2\theta)\right]_0^{\frac\pi2}
 =\frac{\pi a^2}{4}.
\end{aligned}
\]
曲線は $x^2+y^2=a^2$ の上半分であり，積分は半径 $a$ の円の四分の一の面積だから，幾何学的な結果とも一致する。積分計算そのものは，三角置換と二倍角の公式で得られている。

$0\le\theta\le\pi$ を選ぶと，$\sin\theta$ は $0$ から $1$ へ増えた後，$0$ へ戻る。この区間全体を $0\le x\le a$ と一対一に対応させることはできず，後半では $\cos\theta$ も負になる。枝を指定することは，こうした重複と符号を確定するために必要なのである。

もう一つ，$1+x^2$ には $1+\tan^2\theta=\frac1{\cos^2\theta}$ が対応する。$x=\tan\theta$ と置くと，微分も $\dd x=\frac1{\cos^2\theta}\dd\theta$ になる。この一致により，分母が消える場合がある。

\begin{example}{分母と微分の係数が打ち消し合う}
$\int_0^1\frac1{1+x^2}\dd x$ を求めよ。
\end{example}

$x=\tan\theta$ と置き，$0\le\theta\le\frac\pi4$ を選ぶ。この区間で正接は狭義単調増加し，端点は $x=0,1$ に対応する。よって
\[
\begin{aligned}
 \int_0^1\frac1{1+x^2}\dd x
 &=\int_0^{\frac\pi4}
   \frac1{1+\tan^2\theta}\,
   \frac1{\cos^2\theta}\dd\theta\\
 &=\int_0^{\frac\pi4}1\dd\theta
 =\frac\pi4.
\end{aligned}
\]
不定積分として表すなら，最後に $\tan\theta=x$ を満たす $\theta$ を $x$ の関数として書くことになる。この逆関数の名前，枝，微分は\chapterref{章:置換の背景}で詳しく扱う。ここでは定積分の端点が既知の角に対応するため，逆関数の記号を導入しなくても計算を終えられる。

